\documentclass[aspectratio=169]{beamer}

\usetheme{Madrid}
\usepackage[spanish]{babel}
\usepackage[utf8]{inputenc}
\usepackage{textcomp}
\usepackage[T1]{fontenc}
\usepackage{lmodern}
\usepackage{graphicx}
\usepackage{booktabs}
\usepackage{verbatim}
\usepackage{amssymb}
\usepackage{pifont}
\usepackage{moreverb}
\usepackage[most]{tcolorbox}
\usepackage{tikz}
\usetikzlibrary{shapes, arrows, positioning}

% Configuración de TikZ
\tikzstyle{block} = [
  rectangle, draw, rounded corners,
  text centered, minimum height=1cm,
  minimum width=2.5cm, font=\small, fill=gray!5
]
\tikzstyle{arrow} = [->, thick]

% Configuración de tcolorbox
\tcbset{
  colback=gray!5, colframe=gray!50,
  boxrule=0.4pt, sharp corners,
}

\title[BookMate]{BookMate — Sistema de Recomendación de Libros Académicos}
\subtitle{Análisis Conceptual: Especificaciones UC, Modelo de Análisis y Arquitectura C4}
\author[GRUPO 6]{Delgado R., G. \and Osorio M., A. \and Rojas A., J. \and Torres R., J. \and Valverde G., Y. \and Villanueva A., F.}
\institute[FC-UNI]{UNIVERSIDAD NACIONAL DE INGENIERÍA}
\titlegraphic{\includegraphics[height=1.4cm]{example-image}}
% NOTA: Reemplazar 'example-image' con 'Imagenes/logo_uni.png' cuando esté disponible

\date{Octubre 2025}

\begin{document}

% Portada
\begin{frame}
\titlepage
\end{frame}

% Agenda
\begin{frame}
\frametitle{Agenda}
\tableofcontents
\end{frame}

\section{Introducción al Análisis}

\begin{frame}
\frametitle{Fase de Análisis Conceptual}

\begin{columns}
\begin{column}{0.5\textwidth}
\textbf{Objetivo de esta Fase:}
\begin{itemize}
\item Comprender QUÉ debe hacer el sistema
\item Identificar entidades del dominio
\item Definir comportamientos conceptuales
\item Diseñar arquitectura conceptual
\end{itemize}

\vspace{0.3cm}

\textbf{Enfoque:}
\begin{tcolorbox}[colback=blue!5, colframe=blue!50]
\textbf{Sin tecnologías específicas}\\
Nos enfocamos en el comportamiento y la estructura conceptual, no en la implementación técnica.
\end{tcolorbox}
\end{column}

\begin{column}{0.5\textwidth}
\textbf{Entregables de Semana 11:}
\begin{enumerate}
\item Especificaciones detalladas de UC
\item Modelo de Análisis (BCE)
\item Diagramas de Robustez
\item Diagramas de Secuencia
\item Arquitectura C4 (Niveles 1-2)
\item Diagramas de Estados
\item Tarjetas CRC
\end{enumerate}
\end{column}
\end{columns}

\vspace{0.3cm}

\begin{tcolorbox}[colback=yellow!5, colframe=yellow!50]
\textbf{Nota:} Las tecnologías específicas (Spring Boot, PostgreSQL, Python) se definirán en la fase de diseño (Semana 13).
\end{tcolorbox}

\end{frame}

\begin{frame}
\frametitle{Modelo Boundary-Control-Entity (BCE)}

\begin{columns}
\begin{column}{0.33\textwidth}
\textbf{Boundary (Frontera)}
\begin{itemize}
\item Interfaz de Usuario
\item Interfaz de Búsqueda
\item Interfaz de Detalle
\item Interfaz de Admin
\item Interfaz con IA
\end{itemize}
\end{column}

\begin{column}{0.33\textwidth}
\textbf{Control (Controlador)}
\begin{itemize}
\item Gestor de Catálogo
\item Gestor de Búsqueda
\item Generador de Recomendaciones
\item Analizador Semántico
\item Calculador Heurístico
\item Validador de Datos
\end{itemize}
\end{column}

\begin{column}{0.33\textwidth}
\textbf{Entity (Entidad)}
\begin{itemize}
\item Libro
\item Autor
\item Representación Semántica
\item Recomendación
\item Resultado de Búsqueda
\item Usuario
\end{itemize}
\end{column}
\end{columns}

\vspace{0.5cm}

\begin{center}
\begin{tikzpicture}[node distance=2cm]
\node[block, fill=LightBlue] (boundary) {Boundary};
\node[block, fill=LightGreen, right=of boundary] (control) {Control};
\node[block, fill=Pink, right=of control] (entity) {Entity};
\draw[arrow] (boundary) -- (control);
\draw[arrow] (control) -- (entity);
\end{tikzpicture}
\end{center}

\end{frame}

\section{Especificaciones de Casos de Uso}

\begin{frame}
\frametitle{Casos de Uso Priorizados}

\begin{table}
\centering
\small
\begin{tabular}{p{1.5cm}p{4.5cm}p{2cm}p{2cm}}
\toprule
\textbf{ID} & \textbf{Nombre} & \textbf{Prioridad} & \textbf{Complejidad} \\
\midrule
\textbf{UC-06} & Obtener Recomendaciones & \textbf{Muy Alta} ⭐ & Alta \\
UC-01 & Crear Libro & Alta & Media \\
UC-05 & Buscar Libros & Alta & Media \\
UC-03 & Actualizar Libro & Alta & Media \\
UC-04 & Eliminar Libro & Alta & Baja \\
\bottomrule
\end{tabular}
\end{table}

\vspace{0.3cm}

\begin{tcolorbox}[colback=red!5, colframe=red!50]
\textbf{UC-06 es el caso de uso más crítico} - Define el valor principal del sistema: recomendaciones inteligentes con IA.
\end{tcolorbox}

\end{frame}

\begin{frame}
\frametitle{UC-06: Obtener Recomendaciones (Caso Crítico)}

\textbf{Descripción:} El sistema genera una lista de libros similares usando análisis semántico o reglas heurísticas.

\textbf{Actores:}
\begin{itemize}
\item \textbf{Principal:} Usuario
\item \textbf{Secundario:} Componente de Inteligencia Artificial
\end{itemize}

\textbf{Precondiciones:}
\begin{itemize}
\item Libro de referencia existe en el sistema
\item Catálogo tiene al menos 7 libros
\end{itemize}

\textbf{Postcondiciones:}
\begin{itemize}
\item Se retornan 6 libros similares
\item Tiempo de respuesta <3 segundos (p95)
\item Al menos 70\% de relevancia (si usa IA)
\end{itemize}

\vspace{0.3cm}

\begin{tcolorbox}[colback=green!5, colframe=green!50]
\textbf{Valor para el Negocio:} Muy Alto - Es el diferenciador principal del sistema
\end{tcolorbox}

\end{frame}

\begin{frame}
\frametitle{UC-06: Flujo Principal (Análisis Semántico)}

\begin{enumerate}
\item Usuario solicita recomendaciones para libro X
\item Sistema recupera datos del libro X
\item Sistema verifica disponibilidad del componente IA
\item Sistema envía solicitud a IA con sinopsis
\item \textbf{IA genera embedding} de la sinopsis
\item \textbf{IA calcula similitud de coseno} con todos los libros
\item IA retorna top 6 libros más similares
\item Sistema consulta detalles de libros recomendados
\item Sistema muestra recomendaciones al usuario
\end{enumerate}

\vspace{0.3cm}

\begin{center}
\includegraphics[width=0.7\textwidth]{example-image-a}
% NOTA: Reemplazar con diagrama de secuencia del flujo principal de UC-06
% Mostrar: Usuario → Sistema → IA → Sistema → Usuario
% Incluir pasos de generación de embeddings y cálculo de similitud
\end{center}

\end{frame}

\begin{frame}
\frametitle{UC-06: Flujo Alternativo (Heurísticas)}

\textbf{Punto de divergencia:} Componente IA no disponible

\begin{enumerate}
\item Sistema detecta que IA no está disponible
\item Sistema activa modo de respaldo (heurísticas)
\item Sistema aplica algoritmo heurístico:
  \begin{itemize}
  \item Género compartido: +3 puntos
  \item Autor común: +5 puntos
  \item Tags compartidos: +2 puntos cada uno
  \item Precio similar (±20\%): +1 punto
  \item Rating alto (≥4.0): +1 punto
  \end{itemize}
\item Sistema ordena por puntuación descendente
\item Sistema retorna top 6 libros
\end{enumerate}

\vspace{0.3cm}

\begin{tcolorbox}[colback=yellow!5, colframe=yellow!50]
\textbf{Ventaja:} El sistema siempre retorna recomendaciones, garantizando alta disponibilidad.
\end{tcolorbox}

\end{frame}

\begin{frame}
\frametitle{UC-01: Crear Libro - Flujo Principal}

\begin{columns}
\begin{column}{0.5\textwidth}
\textbf{Pasos:}
\begin{enumerate}
\item Admin completa formulario
\item Sistema valida datos
\item Sistema crea libro en BD
\item Sistema notifica a IA
\item IA genera embedding
\item Sistema confirma creación
\end{enumerate}
\end{column}

\begin{column}{0.5\textwidth}
\textbf{Validaciones:}
\begin{itemize}
\item Título no vacío
\item Autor existe
\item Precio > 0
\item Fecha válida
\item ISBN único (si se proporciona)
\end{itemize}
\end{column}
\end{columns}

\vspace{0.3cm}

\begin{center}
\includegraphics[width=0.6\textwidth]{example-image-b}
% NOTA: Reemplazar con diagrama de robustez de UC-01
% Mostrar: Admin → InterfazAdmin → GestorCatálogo → Validador → Libro → AnalizadorSemántico
\end{center}

\end{frame}

\section{Modelo de Análisis}

\begin{frame}
\frametitle{Entidades del Dominio}

\begin{table}
\centering
\tiny
\begin{tabular}{p{3cm}p{4cm}p{4cm}}
\toprule
\textbf{Entidad} & \textbf{Atributos Principales} & \textbf{Responsabilidades} \\
\midrule
\textbf{Libro} & Título, autor, género, precio, sinopsis, ISBN & Mantener información completa, proveer datos para búsqueda \\
\textbf{Autor} & Nombre, biografía, nacionalidad & Mantener información del autor, relacionar libros \\
\textbf{Representación Semántica} & Vector numérico, fecha generación & Almacenar embeddings, proveer datos para similitud \\
\textbf{Recomendación} & Libro recomendado, puntuación, método & Representar resultado de similitud calculada \\
\textbf{Resultado de Búsqueda} & Lista libros, consulta original, relevancia & Agrupar resultados de búsqueda \\
\textbf{Usuario} & Nombre, correo, rol & Mantener información, determinar permisos \\
\bottomrule
\end{tabular}
\end{table}

\end{frame}

\begin{frame}
\frametitle{Controladores (Lógica de Negocio)}

\begin{columns}
\begin{column}{0.5\textwidth}
\textbf{Controladores Principales:}
\begin{itemize}
\item \textbf{Gestor de Catálogo}
  \begin{itemize}
  \item Coordina CRUD de libros/autores
  \item Valida datos
  \item Notifica a IA
  \end{itemize}
\item \textbf{Gestor de Búsqueda}
  \begin{itemize}
  \item Coordina búsquedas
  \item Calcula relevancia
  \item Ordena resultados
  \end{itemize}
\item \textbf{Generador de Recomendaciones}
  \begin{itemize}
  \item Decide método (IA/heurístico)
  \item Coordina generación
  \item Maneja fallback
  \end{itemize}
\end{itemize}
\end{column}

\begin{column}{0.5\textwidth}
\textbf{Controladores Especializados:}
\begin{itemize}
\item \textbf{Analizador Semántico}
  \begin{itemize}
  \item Genera embeddings
  \item Calcula similitud
  \item Comunica con IA
  \end{itemize}
\item \textbf{Calculador Heurístico}
  \begin{itemize}
  \item Aplica reglas
  \item Calcula puntuaciones
  \item Ordena resultados
  \end{itemize}
\item \textbf{Validador de Datos}
  \begin{itemize}
  \item Verifica integridad
  \item Valida formatos
  \item Genera errores
  \end{itemize}
\end{itemize}
\end{column}
\end{columns}

\end{frame}

\begin{frame}
\frametitle{Interfaces (Fronteras)}

\begin{table}
\centering
\small
\begin{tabular}{p{3.5cm}p{6cm}p{2cm}}
\toprule
\textbf{Interfaz} & \textbf{Responsabilidades} & \textbf{Actor} \\
\midrule
Interfaz de Catálogo & Mostrar lista de libros, navegación & Usuario \\
Interfaz de Búsqueda & Capturar consulta, mostrar resultados & Usuario \\
Interfaz de Detalle & Mostrar libro completo, recomendaciones & Usuario \\
Interfaz de Administración & Formularios CRUD, gestión de catálogo & Administrador \\
Interfaz con IA & Comunicar con servicio externo de IA & Sistema IA \\
\bottomrule
\end{tabular}
\end{table}

\vspace{0.3cm}

\begin{center}
\includegraphics[width=0.7\textwidth]{example-image-c}
% NOTA: Reemplazar con diagrama de arquitectura BCE mostrando las 3 capas
% Boundary (arriba) → Control (medio) → Entity (abajo)
% Con ejemplos de cada tipo de elemento
\end{center}

\end{frame}

\section{Diagramas de Robustez}

\begin{frame}
\frametitle{Diagrama de Robustez: UC-06 Recomendaciones}

\begin{center}
\includegraphics[width=0.9\textwidth]{example-image}
% NOTA: Reemplazar con diagrama de robustez completo de UC-06
% Mostrar:
% - Boundary: Interfaz de Detalle, Interfaz con IA
% - Control: Generador de Recomendaciones, Analizador Semántico, Calculador Heurístico
% - Entity: Libro, Catálogo, Representación Semántica, Recomendación
% - Flujos principales y alternativos
\end{center}

\textbf{Elementos clave:}
\begin{itemize}
\item \textbf{Boundary:} Interfaz de Detalle, Interfaz con IA
\item \textbf{Control:} Generador de Recomendaciones, Analizador Semántico, Calculador Heurístico
\item \textbf{Entity:} Libro, Catálogo, Representación Semántica, Recomendación
\end{itemize}

\end{frame}

\begin{frame}
\frametitle{Diagrama de Robustez: UC-01 Crear Libro}

\begin{center}
\includegraphics[width=0.8\textwidth]{example-image-a}
% NOTA: Reemplazar con diagrama de robustez de UC-01
% Mostrar:
% - Admin → InterfazAdmin → GestorCatálogo → Validador → Libro → AnalizadorSemántico
% - Flujo de validación y creación
% - Notificación a IA para generar embedding
\end{center}

\textbf{Flujo:}
\begin{enumerate}
\item Admin → Interfaz de Administración
\item Interfaz → Gestor de Catálogo
\item Gestor → Validador de Datos
\item Gestor → Libro (crea entidad)
\item Gestor → Analizador Semántico (notifica)
\item Analizador → Interfaz con IA
\item Analizador → Representación Semántica (almacena)
\end{enumerate}

\end{frame}

\begin{frame}
\frametitle{Vista General de Robustez del Sistema}

\begin{center}
\includegraphics[width=0.85\textwidth]{example-image-b}
% NOTA: Reemplazar con diagrama de robustez general
% Mostrar todos los elementos BCE organizados:
% - Boundaries en la parte superior
% - Controls en el medio
% - Entities en la parte inferior
% - Relaciones entre ellos
\end{center}

\textbf{Organización:}
\begin{itemize}
\item \textbf{Arriba:} Boundaries (Interfaces con actores)
\item \textbf{Medio:} Controls (Lógica de negocio)
\item \textbf{Abajo:} Entities (Datos del dominio)
\end{itemize}

\end{frame}

\section{Diagramas de Secuencia de Análisis}

\begin{frame}
\frametitle{Diagramas de Secuencia: Concepto}

\begin{columns}
\begin{column}{0.5\textwidth}
\textbf{¿Qué son?}
\begin{itemize}
\item Diagramas de interacción que muestran el orden temporal
\item Colaboración entre elementos a lo largo del tiempo
\item Mensajes intercambiados entre actores, boundaries, controls y entities
\item Flujos alternativos y decisiones
\end{itemize}

\vspace{0.3cm}

\textbf{Propósito:}
\begin{itemize}
\item Visualizar orden temporal de operaciones
\item Entender flujo completo de casos de uso
\item Documentar comportamiento del sistema
\end{itemize}
\end{column}

\begin{column}{0.5\textwidth}
\textbf{Diferencia con Robustez:}
\begin{table}
\centering
\small
\begin{tabular}{p{2.5cm}p{3.5cm}}
\toprule
\textbf{Aspecto} & \textbf{Enfoque} \\
\midrule
\textbf{Robustez} & Estructura estática \\
\textbf{Secuencia} & Orden temporal \\
\bottomrule
\end{tabular}
\end{table}

\vspace{0.3cm}

\textbf{Relación:}
\begin{itemize}
\item Robustez: \textit{¿Qué elementos participan?}
\item Secuencia: \textit{¿En qué orden interactúan?}
\end{itemize}
\end{column}
\end{columns}

\vspace{0.3cm}

\begin{tcolorbox}[colback=blue!5, colframe=blue!50]
\textbf{Nota:} Los diagramas de secuencia de análisis son conceptuales y no mencionan tecnologías específicas. Utilizan la misma notación BCE (Boundary-Control-Entity).
\end{tcolorbox}

\end{frame}

\begin{frame}
\frametitle{Diagrama de Secuencia: UC-06 Recomendaciones}

\begin{center}
\includegraphics[width=0.9\textwidth]{example-image}
% NOTA: Reemplazar con diagrama de secuencia de UC-06
% Mostrar orden temporal:
% 1. Usuario → InterfazDetalle (solicita recomendaciones)
% 2. InterfazDetalle → Generador (solicitud)
% 3. Generador → LibroRef (consulta datos)
% 4. Generador → AnalizadorSem (solicita análisis)
% 5. AnalizadorSem → InterfazIA (envía sinopsis)
% 6. InterfazIA → AnalizadorSem (retorna IDs similares)
% 7. Generador → Catalogo (obtiene detalles)
% 8. Generador → InterfazDetalle (retorna recomendaciones)
% 9. InterfazDetalle → Usuario (muestra resultados)
\end{center}

\textbf{Flujo temporal:}
\begin{enumerate}
\item Usuario solicita recomendaciones
\item Sistema consulta libro de referencia
\item Sistema decide método (IA o heurístico)
\item Sistema calcula similitud
\item Sistema obtiene detalles de libros recomendados
\item Sistema muestra recomendaciones al usuario
\end{enumerate}

\end{frame}

\begin{frame}
\frametitle{Diagrama de Secuencia: UC-06 Flujo Alternativo}

\begin{center}
\includegraphics[width=0.85\textwidth]{example-image-a}
% NOTA: Reemplazar con diagrama de secuencia mostrando flujo alternativo
% Mostrar fallback a heurísticas cuando IA no está disponible
\end{center}

\textbf{Flujo alternativo (Fallback a Heurísticas):}
\begin{enumerate}
\item Sistema detecta que Analizador Semántico no está disponible
\item Sistema activa Calculador Heurístico
\item Calculador aplica reglas de similitud:
  \begin{itemize}
  \item Género compartido: +3 puntos
  \item Autor común: +5 puntos
  \item Tags compartidos: +2 puntos cada uno
  \item Precio similar: +1 punto
  \item Rating alto: +1 punto
  \end{itemize}
\item Sistema ordena por puntuación
\item Sistema retorna top 6 libros heurísticos
\end{enumerate}

\vspace{0.3cm}

\begin{tcolorbox}[colback=green!5, colframe=green!50]
\textbf{Ventaja:} El sistema siempre retorna recomendaciones, garantizando alta disponibilidad mediante fallback automático.
\end{tcolorbox}

\end{frame}

\begin{frame}
\frametitle{Diagrama de Secuencia: UC-01 Crear Libro}

\begin{center}
\includegraphics[width=0.8\textwidth]{example-image-b}
% NOTA: Reemplazar con diagrama de secuencia de UC-01
% Mostrar orden temporal de creación de libro
\end{center}

\textbf{Secuencia temporal:}
\begin{enumerate}
\item Administrador completa formulario
\item Interfaz envía datos a Gestor de Catálogo
\item Gestor solicita validación
\item Validador verifica datos y autor
\item Gestor crea entidad Libro
\item Gestor notifica a Analizador Semántico
\item Analizador genera representación semántica
\item Sistema confirma creación
\end{enumerate}

\vspace{0.3cm}

\textbf{Flujo alternativo:} Si datos son inválidos, sistema retorna errores específicos sin crear el libro.

\end{frame}

\begin{frame}
\frametitle{Diagrama de Secuencia: UC-05 Buscar Libros}

\begin{center}
\includegraphics[width=0.8\textwidth]{example-image-c}
% NOTA: Reemplazar con diagrama de secuencia de UC-05
% Mostrar orden temporal de búsqueda
\end{center}

\textbf{Secuencia temporal:}
\begin{enumerate}
\item Usuario ingresa consulta de búsqueda
\item Interfaz envía solicitud a Gestor de Búsqueda
\item Gestor busca en Catálogo (título, género)
\item Gestor busca en Autores (nombre)
\item Gestor calcula relevancia por tipo de coincidencia
\item Gestor ordena resultados por relevancia
\item Gestor crea entidad Resultado de Búsqueda
\item Interfaz muestra resultados al usuario
\end{enumerate}

\vspace{0.3cm}

\textbf{Flujo alternativo:} Si no hay resultados, sistema muestra libros populares como sugerencias.

\end{frame}

\begin{frame}
\frametitle{Robustez vs Secuencia: Complementariedad}

\begin{columns}
\begin{column}{0.5\textwidth}
\textbf{Diagrama de Robustez}
\begin{itemize}
\item \textbf{Enfoque:} Estructura estática
\item \textbf{Muestra:} ¿Qué elementos participan?
\item \textbf{Propósito:} Validar completitud
\item \textbf{Vista:} Relaciones entre elementos
\end{itemize}

\vspace{0.3cm}

\textbf{Ejemplo:}
\begin{itemize}
\item Boundary: Interfaz de Detalle
\item Control: Generador de Recomendaciones
\item Entity: Libro, Recomendación
\end{itemize}
\end{column}

\begin{column}{0.5\textwidth}
\textbf{Diagrama de Secuencia}
\begin{itemize}
\item \textbf{Enfoque:} Orden temporal
\item \textbf{Muestra:} ¿En qué orden interactúan?
\item \textbf{Propósito:} Documentar comportamiento
\item \textbf{Vista:} Secuencia de mensajes
\end{itemize}

\vspace{0.3cm}

\textbf{Ejemplo:}
\begin{enumerate}
\item Usuario → Interfaz
\item Interfaz → Generador
\item Generador → Libro
\item Generador → Analizador
\item ...
\end{enumerate}
\end{column}
\end{columns}

\vspace{0.5cm}

\begin{tcolorbox}[colback=yellow!5, colframe=yellow!50]
\textbf{Conclusión:} Ambos diagramas son necesarios. El diagrama de robustez valida la estructura, mientras que el diagrama de secuencia documenta el comportamiento temporal.
\end{tcolorbox}

\end{frame}

\section{Arquitectura C4}

\begin{frame}
\frametitle{C4 Model - Nivel 1: Contexto}

\begin{center}
\includegraphics[width=0.9\textwidth]{example-image-c}
% NOTA: Reemplazar con diagrama C4 Nivel 1 (Contexto)
% Mostrar:
% - Personas: Usuario, Administrador
% - Sistema BookMate (centro)
% - Sistema de IA (externo)
% - Relaciones entre ellos
\end{center}

\textbf{Elementos:}
\begin{itemize}
\item \textbf{Usuarios:} Estudiante, profesor, investigador
\item \textbf{Administradores:} Bibliotecario, gestor de catálogo
\item \textbf{Sistema BookMate:} Sistema principal
\item \textbf{Sistema de IA:} Servicio externo para análisis semántico
\end{itemize}

\end{frame}

\begin{frame}
\frametitle{C4 Model - Nivel 2: Contenedores}

\begin{center}
\includegraphics[width=0.85\textwidth]{example-image}
% NOTA: Reemplazar con diagrama C4 Nivel 2 (Contenedores)
% Mostrar:
% - Aplicación Web (Frontend)
% - Sistema de Gestión (Backend)
% - Sistema de Almacenamiento (BD)
% - Servicio de IA (externo)
% - Flujos de comunicación entre contenedores
\end{center}

\textbf{Contenedores:}
\begin{itemize}
\item \textbf{Aplicación Web:} Interfaz de usuario (HTML/CSS/JS)
\item \textbf{Sistema de Gestión:} Lógica de negocio (Backend)
\item \textbf{Sistema de Almacenamiento:} Base de datos
\item \textbf{Servicio de IA:} Microservicio externo (Python)
\end{itemize}

\end{frame}

\begin{frame}
\frametitle{Flujo Entre Contenedores: Recomendaciones}

\begin{center}
\includegraphics[width=0.8\textwidth]{example-image-a}
% NOTA: Reemplazar con diagrama de flujo entre contenedores para UC-06
% Mostrar secuencia:
% 1. Usuario → WebApp
% 2. WebApp → Backend
% 3. Backend → BD (consulta libro)
% 4. Backend → Servicio IA
% 5. Servicio IA → BD (lee embeddings)
% 6. Servicio IA → Backend (retorna IDs)
% 7. Backend → BD (consulta detalles)
% 8. Backend → WebApp
% 9. WebApp → Usuario
\end{center}

\textbf{Flujo:}
\begin{enumerate}
\item Usuario solicita recomendaciones
\item Backend consulta BD
\item Backend solicita a Servicio IA
\item Servicio IA calcula similitud
\item Backend retorna resultados
\end{enumerate}

\end{frame}

\begin{frame}
\frametitle{Descomposición del Backend (Componentes)}

\begin{center}
\includegraphics[width=0.9\textwidth]{example-image-b}
% NOTA: Reemplazar con diagrama C4 mostrando componentes internos del backend
% Mostrar:
% - Gestor de Catálogo
% - Gestor de Búsqueda
% - Generador de Recomendaciones
% - Analizador Semántico
% - Calculador Heurístico
% - Validador de Datos
% - Relaciones entre componentes
\end{center}

\textbf{Componentes Conceptuales:}
\begin{itemize}
\item Gestor de Catálogo, Gestor de Búsqueda
\item Generador de Recomendaciones
\item Analizador Semántico, Calculador Heurístico
\item Validador de Datos
\end{itemize}

\end{frame}

\section{Diagramas de Estados}

\begin{frame}
\frametitle{Estado: Ciclo de Vida de un Libro}

\begin{center}
\includegraphics[width=0.8\textwidth]{example-image-c}
% NOTA: Reemplazar con diagrama de estados del ciclo de vida de un libro
% Estados: Borrador → Validando → Registrado → GenerandoEmbedding → Activo
% Transiciones: Editar, Eliminar, etc.
\end{center}

\textbf{Estados principales:}
\begin{itemize}
\item \textbf{Borrador:} Datos ingresados, no validados
\item \textbf{Validando:} Sistema valida datos
\item \textbf{Registrado:} Libro persistido en BD
\item \textbf{GenerandoEmbedding:} Procesando con IA
\item \textbf{Activo:} Disponible para búsquedas y recomendaciones
\end{itemize}

\end{frame}

\begin{frame}
\frametitle{Estado: Solicitud de Recomendación}

\begin{center}
\includegraphics[width=0.75\textwidth]{example-image}
% NOTA: Reemplazar con diagrama de estados de una solicitud de recomendación
% Estados: Solicitada → VerificandoIA → ProcesandoConIA/ProcesandoHeuristico → Completada → Entregada
\end{center}

\textbf{Estados:}
\begin{enumerate}
\item \textbf{Solicitada:} Usuario pide recomendaciones
\item \textbf{VerificandoIA:} Verificando disponibilidad
\item \textbf{ProcesandoConIA:} Análisis semántico (método primario)
\item \textbf{ProcesandoHeuristico:} Reglas de negocio (fallback)
\item \textbf{Completada:} Recomendaciones generadas
\item \textbf{Entregada:} Mostradas al usuario
\end{enumerate}

\end{frame}

\begin{frame}
\frametitle{Estado: Representación Semántica (Embedding)}

\begin{center}
\includegraphics[width=0.7\textwidth]{example-image-a}
% NOTA: Reemplazar con diagrama de estados de una representación semántica
% Estados: Pendiente → Generando → Almacenada → Actualizando → Error
\end{center}

\textbf{Estados:}
\begin{itemize}
\item \textbf{Pendiente:} Libro creado, esperando generación
\item \textbf{Generando:} Procesando sinopsis con NLP
\item \textbf{Almacenada:} Embedding disponible y válido
\item \textbf{Actualizando:} Regenerando (si sinopsis cambió)
\item \textbf{Error:} Fallo en generación (reintento programado)
\end{itemize}

\end{frame}

\section{Tarjetas CRC}

\begin{frame}
\frametitle{Tarjetas CRC: Concepto}

\textbf{CRC = Class-Responsibility-Collaboration}

\begin{columns}
\begin{column}{0.5\textwidth}
\textbf{Class (Clase):}
\begin{itemize}
\item Nombre de la clase conceptual
\item Entidad, Controlador o Frontera
\end{itemize}

\textbf{Responsibilities:}
\begin{itemize}
\item Qué hace la clase
\item Funciones principales
\item Responsabilidades específicas
\end{itemize}
\end{column}

\begin{column}{0.5\textwidth}
\textbf{Collaborators:}
\begin{itemize}
\item Con quiénes interactúa
\item Dependencias
\item Relaciones
\end{itemize}
\end{column}
\end{columns}

\vspace{0.5cm}

\begin{center}
\includegraphics[width=0.6\textwidth]{example-image-b}
% NOTA: Reemplazar con ejemplo visual de una tarjeta CRC
% Mostrar formato: Clase arriba, Responsabilidades a la izquierda, Colaboradores a la derecha
\end{center}

\end{frame}

\begin{frame}
\frametitle{Tarjetas CRC: Ejemplos Clave}

\textbf{CRC-09: GeneradorDeRecomendaciones}

\begin{table}
\centering
\small
\begin{tabular}{p{5cm}p{7cm}}
\toprule
\textbf{Responsabilidades} & \textbf{Colaboradores} \\
\midrule
• Coordinar generación de recomendaciones & • Libro \\
• Determinar método (IA/heurístico) & • AnalizadorSemántico \\
• Verificar disponibilidad de IA & • CalculadorHeurístico \\
• Activar método apropiado & • Recomendación \\
• Manejar fallback transparente & • InterfazDeDetalle \\
\bottomrule
\end{tabular}
\end{table}

\vspace{0.3cm}

\textbf{CRC-10: AnalizadorSemántico}

\begin{table}
\centering
\small
\begin{tabular}{p{5cm}p{7cm}}
\toprule
\textbf{Responsabilidades} & \textbf{Colaboradores} \\
\midrule
• Generar representación vectorial & • Libro \\
• Calcular similitud entre embeddings & • RepresentaciónSemántica \\
• Comunicarse con componente IA & • InterfazConIA \\
• Retornar libros más similares & • GeneradorDeRecomendaciones \\
\bottomrule
\end{tabular}
\end{table}

\end{frame}

\begin{frame}
\frametitle{Resumen de Tarjetas CRC Generadas}

\begin{table}
\centering
\tiny
\begin{tabular}{p{2.5cm}p{2cm}p{2cm}p{2cm}}
\toprule
\textbf{Categoría} & \textbf{Cantidad} & \textbf{Ejemplos} & \textbf{Prioridad} \\
\midrule
\textbf{Entidades} & 6 & Libro, Autor, RepresentaciónSemántica & Alta \\
\textbf{Controladores} & 6 & GestorCatálogo, GeneradorRecomendaciones & Muy Alta \\
\textbf{Interfaces} & 5 & InterfazDetalle, InterfazAdmin & Alta \\
\textbf{Auxiliares} & 3 & Catálogo, ConjuntoRecomendaciones & Media \\
\bottomrule
\end{tabular}
\end{table}

\vspace{0.3cm}

\textbf{Total: 20 tarjetas CRC}

\begin{tcolorbox}[colback=green!5, colframe=green!50]
Las tarjetas CRC proporcionan una base sólida para el diseño detallado en Semana 15, donde se transformarán en clases de implementación con métodos específicos.
\end{tcolorbox}

\end{frame}

\section{Flujos de Información}

\begin{frame}
\frametitle{Flujo: Obtener Recomendaciones (Detallado)}

\begin{center}
\includegraphics[width=0.9\textwidth]{example-image-c}
% NOTA: Reemplazar con diagrama de flujo detallado de recomendaciones
% Mostrar secuencia completa:
% Usuario → InterfazDetalle → GeneradorRecomendaciones → AnalizadorSemántico → InterfazIA → SistemaIA
% Y retorno con resultados
\end{center}

\textbf{Secuencia:}
\begin{enumerate}
\item Usuario → Interfaz de Detalle
\item Interfaz → Generador de Recomendaciones
\item Generador → Analizador Semántico
\item Analizador → Interfaz con IA
\item Interfaz → Sistema de IA (externo)
\item Sistema IA procesa y retorna
\item Analizador → Generador
\item Generador → Interfaz
\item Interfaz → Usuario
\end{enumerate}

\end{frame}

\begin{frame}
\frametitle{Flujo: Crear Libro (Detallado)}

\begin{center}
\includegraphics[width=0.85\textwidth]{example-image}
% NOTA: Reemplazar con diagrama de flujo de creación de libro
% Mostrar: Admin → InterfazAdmin → GestorCatálogo → Validador → Libro → AnalizadorSemántico → IA
\end{center}

\textbf{Secuencia:}
\begin{enumerate}
\item Admin → Interfaz de Administración
\item Interfaz → Gestor de Catálogo
\item Gestor → Validador de Datos
\item Validador valida → Gestor
\item Gestor → Libro (crea entidad)
\item Gestor → Analizador Semántico
\item Analizador → Interfaz con IA
\item IA genera embedding
\item Analizador → Representación Semántica (almacena)
\end{enumerate}

\end{frame}

\section{Resumen y Conclusiones}

\begin{frame}
\frametitle{Resumen del Análisis Conceptual}

\begin{columns}
\begin{column}{0.5\textwidth}
\textbf{Entidades Identificadas:}
\begin{itemize}
\item 6 entidades del dominio
\item Libro, Autor, RepresentaciónSemántica
\item Recomendación, ResultadoBúsqueda, Usuario
\end{itemize}

\textbf{Controladores Identificados:}
\begin{itemize}
\item 6 controladores de lógica
\item GestorCatálogo, GestorBúsqueda
\item GeneradorRecomendaciones
\item AnalizadorSemántico, CalculadorHeurístico
\end{itemize}
\end{column}

\begin{column}{0.5\textwidth}
\textbf{Interfaces Identificadas:}
\begin{itemize}
\item 5 interfaces de frontera
\item InterfazCatálogo, InterfazBúsqueda
\item InterfazDetalle, InterfazAdmin
\item InterfazConIA
\end{itemize}

\textbf{Diagramas Generados:}
\begin{itemize}
\item 5 diagramas de robustez
\item 3 diagramas de secuencia de análisis
\item 2 niveles C4 (Contexto, Contenedores)
\item 4 diagramas de estados
\item 20 tarjetas CRC
\end{itemize}
\end{column}
\end{columns}

\end{frame}

\begin{frame}
\frametitle{Arquitectura Conceptual Identificada}

\begin{center}
\includegraphics[width=0.8\textwidth]{example-image-a}
% NOTA: Reemplazar con diagrama de arquitectura conceptual completa
% Mostrar la estructura BCE completa con todos los elementos identificados
\end{center}

\textbf{Organización:}
\begin{itemize}
\item \textbf{Capa de Frontera:} 5 interfaces para interacción con actores
\item \textbf{Capa de Control:} 6 controladores coordinando comportamientos
\item \textbf{Capa de Entidad:} 6 entidades representando el dominio
\end{itemize}

\end{frame}

\begin{frame}
\frametitle{Próximos Pasos: Fase de Diseño}

\textbf{Semana 13 - Diseño Arquitectónico:}

\begin{itemize}
\item Definir stack tecnológico específico
  \begin{itemize}
  \item Backend: Spring Boot
  \item Base de datos: PostgreSQL
  \item Servicio IA: Python + Sentence Transformers
  \item Frontend: HTML/CSS/JavaScript
  \end{itemize}
\item C4 Niveles 3-4 (Componentes y Código)
\item Patrones de diseño aplicados
\item Decisiones arquitectónicas clave
\item Plan de desarrollo del software
\end{itemize}

\vspace{0.3cm}

\begin{tcolorbox}[colback=blue!5, colframe=blue!50]
\textbf{En la fase de diseño:} Transformaremos estos elementos conceptuales en clases de implementación con tecnologías específicas, métodos reales y arquitectura física detallada.
\end{tcolorbox}

\end{frame}

\begin{frame}
\frametitle{Logros del Análisis Conceptual}

\begin{enumerate}
\item ✅ \textbf{Especificaciones completas} de 5 casos de uso principales
\item ✅ \textbf{Modelo de análisis BCE} con 17 elementos identificados
\item ✅ \textbf{Diagramas de robustez} validando casos de uso
\item ✅ \textbf{Diagramas de secuencia} documentando orden temporal
\item ✅ \textbf{Arquitectura C4} niveles 1-2 documentada
\item ✅ \textbf{Diagramas de estados} para entidades clave
\item ✅ \textbf{20 tarjetas CRC} con responsabilidades claras
\item ✅ \textbf{Flujos de información} documentados
\item ✅ \textbf{Separación clara} entre análisis y diseño
\end{enumerate}

\vspace{0.5cm}

\begin{tcolorbox}[colback=green!5, colframe=green!50, title=Objetivo Cumplido]
Hemos comprendido QUÉ debe hacer el sistema y CÓMO se organiza conceptualmente, sin atarnos a tecnologías específicas. Esto proporciona una base sólida para el diseño detallado.
\end{tcolorbox}

\end{frame}

\begin{frame}
\frametitle{¡Preguntas?}

\begin{center}
\Large ¡Gracias por su atención! \\

\vspace{1cm}

\textbf{¿Preguntas sobre el análisis conceptual?}

\vspace{1cm}

\normalsize
\textbf{Grupo 6}\\
Universidad Nacional de Ingeniería\\
CC341 - Ingeniería de Software\\
Ciclo 2025-2

\vspace{0.5cm}

\textbf{Próxima entrega:} Semana 13 - Diseño Arquitectónico (con tecnologías específicas)

\end{center}

\end{frame}

\end{document}

